\section*{Performance Analysis of the Anderson-Accelerated Block Gauss-Seidel Solver}
\label{sec:fp-aa-performance}

\paragraph{Benchmark design and metrics.}
We assess robustness and computational efficiency of the nonlinear coupling solver (block Gauss--Seidel with optional Anderson acceleration; \S\ref{sec:fp-aa}) on the box benchmark under a parabolic compression load.
To isolate coupling behavior, we fix $\Delta t=25~\mathrm{d}$ and $T=500~\mathrm{d}$ (20 steps).
This intentionally large $\Delta t$ is used as a stress test: it is not aimed at resolving stimulus transients ($\tau_S\approx 10~\mathrm{d}$), but at probing the nonlinear contractivity of the coupled map.

We compare two parameter regimes:
\textbf{PhysioSetup} (well-conditioned, more regularized) and \textbf{StiffSetup} (more challenging coupling via reduced stimulus/density diffusion and faster density kinetics).
For each setup we run (i) a pure damped Picard iteration (no Anderson) and (ii) Anderson-accelerated mixing with the parameters in Table~\ref{tab:anderson-params}.

The primary convergence indicator is the coupling residual $r_{\mathrm{Pic}}$, defined by \eqref{eq:fp-picard-res} and evaluated on the coupled state $(\mathbf{L},S,\rho)$.
For diagnostics we also use the per-iteration contraction estimate $\rev{q}_k:=r_k/r_{k-1}$ and the condition number of the regularized Gram matrix (see \S\ref{sec:fp-aa}).
Performance is reported as the cumulative linear-solver time and iteration counts aggregated over the four blocks (mechanics, fabric, stimulus, density).

\begin{table}[htbp]
  \centering
  \caption{Fixed-point and Anderson-acceleration parameters used in the box benchmarks and sweep studies.}
  \label{tab:anderson-params}
  \begin{tabular}{llr}
    \toprule
    \textbf{Parameter} & \textbf{Description} & \textbf{Value} \\
    \midrule
    $m$ & Anderson history size & 5 \\
    $\beta$ & Mixing parameter & 1.0 \\
    $\lambda$ & Relative Tikhonov regularization & $10^{-3}$ \\
    $\mathrm{tol}$ & Coupling tolerance & $10^{-6}$ \\
    $\gamma_{\mathrm{stall}}$ & Restart threshold (stall factor) & 1.1 \\
    $\kappa_{\max}$ & Restart threshold (conditioning) & $10^{4}$ \\
    $c_{\mathrm{step}}$ & Step limiting factor & 1.5 \\
    $w_{\mathrm{stall}}$ & Stall window & 2 \\
    $P_{\mathrm{stall}}$ & Stall patience & 2 \\
    $K_{\max}$ & Max.\ subiterations & 25 \\
    $W$ & Outer-stall window & 4 \\
    $\delta_{\min}$ & Outer-stall minimum relative drop & 0.05 \\
    $P_{\mathrm{outer}}$ & Outer-stall patience & 2 \\
    $\rev{q}_{\mathrm{off}}$ & Anderson deactivation threshold & 0.3 \\
    $\rev{q}_{\mathrm{on}}$ & Anderson activation threshold & 0.7 \\
    $P_{\rev{q}}$ & Anderson switch patience & 4 \\
    \bottomrule
  \end{tabular}
\end{table}

  \begin{figure}[htbp]
  \centering
  \safeincludegraphics[width=1\linewidth]{anderson_comparison.png}
  \caption{Physio vs.\ stiff coupling in the box benchmark: Anderson acceleration vs.\ pure damped Picard iteration.
  (a) Final coupling residual $r_{\mathrm{Pic}}$ per time step (log scale; values below $10^{-14}$ are floored for visibility), with the dashed line marking $\mathrm{tol}=10^{-6}$.
  (b) Total linear-solver time over the full run, stacked by fixed-point termination reason (converged to tolerance, reached the subiteration limit, insufficient progress).
  (c) Representative residual histories: for each configuration we plot the step that required the most subiterations.}
  \label{fig:anderson_comparison}
\end{figure}
\paragraph{Observed robustness trends.}
Figure~\ref{fig:anderson_comparison} highlights two qualitatively different regimes.
In \textbf{PhysioSetup}, the coupling is sufficiently contractive that both Picard and Anderson-accelerated iterations converge rapidly to $\mathrm{tol}$, and the total solve times are similar.
In \textbf{StiffSetup}, the reduced regularization and faster kinetics substantially weaken contractivity, and pure Picard typically fails to satisfy the coupling tolerance within $K_{\max}$ subiterations (termination by subiteration limit or insufficient progress), wasting substantial linear-solver work while leaving a large final residual.
Anderson mixing restores practical robustness by stabilizing the fixed-point iteration and driving the residual to the prescribed tolerance in the same configuration.
\paragraph{Event-driven safeguards and diagnostic run.}
We augment classical Anderson (Pulay/DIIS) mixing with safeguards that are essential in large-scale, strongly coupled runs:
adaptive restarts under stall or ill-conditioning, a step limiter that caps the Anderson update relative to the underlying Picard step, and a hysteresis-based switch that temporarily disables Anderson when the observed contraction is already strong (\S\ref{sec:fp-aa}).
Figure~\ref{fig:anderson_diagnostic} summarizes these mechanisms on a representative StiffSetup run via (i) per-step convergence curves, (ii) global timelines of the residual and contraction, (iii) evolution of the effective history size, (iv) Gram-matrix conditioning, and (v) cumulative event counts (restarts, limiting, Picard-mode switches, and step-level stop reasons).
  
\begin{figure}[htbp]
  \centering
  \safeincludegraphics[width=1\linewidth]{anderson_diagnostic.png}
  \caption{Diagnostic view of the Anderson-accelerated fixed-point solver on a representative stiff run.
  (a) Residual histories $r_{\mathrm{Pic}}$ per step (colored by step index).
  (b) Global residual timeline with an overlaid contraction estimate $\rev{q}_k=r_k/r_{k-1}$.
  (c) Time-step and subiteration counts.
  (d) Effective Anderson history size.
  (e) Conditioning $\kappa(H)$ of the regularized Gram matrix with the restart threshold.
  (f) Cumulative events: stall/conditioning restarts, step limiting, Picard-mode switches, and step-level termination reasons (subiteration limit, insufficient progress).}
  \label{fig:anderson_diagnostic}
\end{figure}

\paragraph{Sensitivity to $(m,\beta,\lambda)$.}
Finally, we quantify parameter sensitivity by a compact $3\times 3\times 3$ sweep over the history size $m$, mixing $\beta$, and relative regularization $\lambda$ (StiffSetup, box benchmark; fixed $\Delta t=25~\mathrm{d}$, $T=100~\mathrm{d}$).
We summarize the mean iterations/step, the contraction estimator, and the Gram conditioning in \cref{fig:anderson_sweep}, highlighting the expected trade-offs: larger $m$ can reduce iterations but tends to deteriorate conditioning, while larger $\lambda$ stabilizes the coefficient solve (and triggers fewer conditioning-related restarts) at the cost of weaker acceleration.

\begin{figure}[htbp]
  \centering
  \safeincludegraphics[width=1\linewidth]{anderson_sweep.png}
  \caption{Sensitivity of coupling performance to the Anderson parameters $(m,\beta,\lambda)$ in the stiff box benchmark (parameter sweep).
  (a) Mean subiterations per step as a function of history size $m$ and mixing $\beta$ at the best-performing $\lambda$ among the sampled values; the red box marks the best cell in the sampled grid.
  (b) Marginal distributions of iterations/step when varying $m$, $\beta$, or $\lambda$ (other parameters averaged out).
  (c) Median contraction $\rev{q}$ as a function of $(m,\beta)$ at the same $\lambda$ as in (a) (smaller is better).
  (d) Median Gram conditioning $\kappa(H)$ versus $m$ for each $\lambda$, with the restart threshold $\kappa_{\max}$ indicated.}
  \label{fig:anderson_sweep}
\end{figure}
